% chapters/01/rangkuman.tex
\bagian{Rangkuman}

\begin{rangkumanbox}
\begin{itemize}
  \item Memprogram permainan berarti menulis program yang menerima masukan,
        memperbarui keadaan, dan menampilkan hasilnya secara terus-menerus;
        karena itu pengelolaan waktu dan penanganan peristiwa menjadi unsur
        pokok, bukan tambahan.
  \item Mata kuliah ini berbobot \BobotSKS\ dengan prasyarat
        \MataKuliahPrasyarat, dan menuntut empat CPMK yang masing-masing
        dibuktikan melalui artefak nyata: prototipe Pygame, mini proyek AI,
        level Unity, dan prototipe Unity lengkap.
  \item Buku ini mengikuti RPS berbasis OBE dengan prinsip penyelarasan
        konstruktif: capaian ditetapkan lebih dahulu, lalu materi, aktivitas,
        dan asesmen diarahkan padanya. Satu bab inti menopang satu Sub-CPMK.
  \item Nomor rujukan pustaka pada buku ini sama dengan nomor rujukan pada
        RPS, yaitu \mbox{[R1]} sampai \mbox{[R35]}.
  \item Materi terbagi menjadi empat blok. Blok I dan Blok II memakai Python
        dan Pygame agar mekanisme permainan terlihat dan algoritma dipahami
        secara mendalam; Blok III dan Blok IV memakai Unity untuk membangun
        level dan mengintegrasikan aset.
  \item Penilaian bersifat autentik dan berbasis kriteria. Syarat kelulusan
        adalah nilai akhir minimal 60, kehadiran minimal 75\%, dan penyerahan
        seluruh keluaran wajib.
  \item Hasil belajar lebih ditentukan oleh kebiasaan kerja mingguan:
        mengetik ulang kode, bereksperimen dengan mengubah nilai, membaca
        pesan kesalahan, mencatat kendala, dan bekerja rutin dalam porsi
        kecil.
\end{itemize}
\end{rangkumanbox}
